% Generated draft: requires mathematical and source review.
\documentclass[11pt]{article}
\usepackage[margin=25mm]{geometry}
\usepackage{fontspec}
\setmainfont{Latin Modern Roman}
\usepackage{amsmath,amssymb,mathtools,unicode-math}
\setmathfont{Latin Modern Math}
\usepackage{enumitem,needspace,xurl,hyperref}
\setlength{\parindent}{0pt}
\setlength{\parskip}{5pt}
\setlength{\emergencystretch}{3em}
\newcommand{\N}{\mathbb{N}}
\newcommand{\Z}{\mathbb{Z}}
\newcommand{\Q}{\mathbb{Q}}
\newcommand{\R}{\mathbb{R}}
\newcommand{\C}{\mathbb{C}}
\newcommand{\F}{\mathbb{F}}
\newcommand{\A}{\mathbb{A}}
\newcommand{\PP}{\mathbb{P}}
\newcommand{\E}{\mathbb{E}}
\newcommand{\eps}{\varepsilon}
\newcommand{\dd}{\,\mathrm{d}}
\begin{document}
\section*{UnsolvedMath ID: 30002883}
\textbf{Source problem code:} catalogue entry
\subsection{Definitions and mathematical statement}
The supplied record contains two related conjectural assertions. Its original wording restricts the first to the ``relevant quantum cluster algebras,'' but does not specify that class. The following formulations are therefore conditional on supplying the indicated algebraic data; a universal assertion about arbitrary quantum cluster algebras is not justified by the record.

For a quantum cluster algebra \(A\) with fixed coefficient and normalization conventions, let \(M_t(a)\) denote the normalized quantum-torus monomial with exponent vector \(a\) in seed \(t\). Let \(D_t\) be the allowed exponent vectors for cluster monomials under those conventions, with nonnegative integer exponents on mutable variables. The treatment of frozen variables and coefficient factors must be specified. Write
\[
\operatorname{CM}(A)=\{M_t(a): t\text{ is a seed of }A,\ a\in D_t\}.
\]
This definition uses normalized quantum monomials, so their equality with basis elements depends on the chosen normalization.

Suppose that a specified realization \(\iota:A\hookrightarrow V\) places \(A\) in an algebra \(V\) carrying a specified dual canonical basis \(\mathcal B^{\vee}\). The first intended assertion is
\[
\forall m\in\operatorname{CM}(A),\qquad
\iota(m)\in\mathcal B^{\vee}.
\]
The question asks whether this holds for every algebra and realization in the intended class. Neither that class nor the ambient algebra, realization, coefficient ring, or dual canonical basis is identified in the supplied record. An arbitrary quantum cluster algebra does not come with such data merely by being a quantum cluster algebra. Literal membership, rather than membership up to a power of the quantum parameter, requires a compatible normalization that is also unspecified.

For the second assertion, let \(\mathcal C\) be a specified tensor category of representations of a quantum affine algebra. Its Grothendieck ring \(K_0(\mathcal C)\) has classes \([X]\), relations
\[
[Y]=[X]+[Z]\quad\text{for each short exact sequence }0\longrightarrow X\longrightarrow Y\longrightarrow Z\longrightarrow0,
\]
and multiplication \([X][Y]=[X\otimes Y]\), whenever these structures are defined in the chosen category. Assume that the classes of simple objects form the intended basis
\[
\mathcal B_{\mathrm{simp}}=\{[L]:L\text{ is simple in }\mathcal C\},
\]
with representatives taken up to isomorphism. Let \(A_{\mathcal C}\) be the relevant cluster algebra with its specified realization \(\jmath:A_{\mathcal C}\hookrightarrow K_0(\mathcal C)\). Here cluster monomials are products of variables from a single cluster, with the normalization and coefficient conventions of that cluster structure. The second intended assertion is
\[
\forall m\in\operatorname{CM}(A_{\mathcal C}),\quad
\exists L\text{ simple in }\mathcal C\quad
\jmath(m)=[L].
\]
The quantum affine algebra, representation category, cluster realization, and conventions ensuring this literal equality are not supplied. These missing choices prevent either assertion from being a fully specified conjecture. No equivalence between the two assertions is asserted here.
\subsection{Short English statement}
The intended questions ask whether every cluster monomial is a dual canonical basis element in the relevant quantum cluster realizations, and whether every cluster monomial in the relevant Grothendieck rings is the class of a simple representation. The relevant families, realizations, coefficients, and normalizations remain unspecified. The supplied record labels the problem partially solved; that label and its dated literature assessment are supplied curation, not independently verified conclusions.
\subsection{Sources}
\begin{itemize}
\item Local UnsolvedMath registry, permanent ID \(30002883\), title ``Cluster Monomials in Dual Canonical Bases,'' exact source problem code \texttt{catalogue entry}, in the Oberwolfach Reports Open Problems collection. The record gives proposed year \(2015\) and names no individual proposer.
\item Supplied report citation: \emph{Enveloping Algebras and Geometric Representation Theory} (2016), Oberwolfach Reports, pp.~1385--1447, DOI \(10.4171/owr/2015/25\). \url{https://doi.org/10.4171/owr/2015/25}. The supplied background attributes the problem to workshop open-problem discussions documented in Oberwolfach Reports. The relationship between the citation year 2016, proposed year 2015, and DOI year component is not resolved here.
\item Kang--Kashiwara--Kim--Oh, \emph{Monoidal categorification}. Supplied literature reference: \url{https://doi.org/10.1090/jams/895}.
\item Gross--Hacking--Keel--Kontsevich, \emph{Canonical bases}. Supplied literature reference: \url{https://arxiv.org/abs/1106.4971}.
\item Literature assessment dated exactly \texttt{2026-08-21}: the supplied curation labels the problem partially solved and reports that cluster-monomial and dual-canonical-basis assertions are proved in substantial monoidal-categorification classes, while the broad conjectural formulation remains unsettled in full generality. This is supplied dated curation derived from the OWR record, not an independently verified literature assessment. The source itself requests independent verification of time-sensitive status and cited claims.
\end{itemize}
\end{document}
